\clearpage\corrigesChapitre{15}{Période et perception sonore}
\corrigesSouples
\begin{blocLivre}{cor-15-1}\begin{corrige}{1}{De la période à la fréquence}
\begin{questions}
\item $T=\qty{2,00e-3}{s}$.
\item $f=1/T=\qty{500}{Hz}$.
\end{questions}
\end{corrige}\end{blocLivre}
\begin{blocLivre}{cor-15-2}\begin{corrige}{2}{De la fréquence à la période}
\begin{questions}
\item $T=1/250=\qty{0,00400}{s}=\qty{4,00}{ms}$.
\item 250 motifs en une seconde, par définition du hertz.
\end{questions}
\end{corrige}\end{blocLivre}
\begin{blocLivre}{cor-15-3}\begin{corrige}{3}{Trois domaines}
\begin{questions}
\item 12 Hz : infrason ; 1000 Hz : domaine audible ; 40 kHz : ultrason.
\item Non. Fréquence et intensité décrivent deux propriétés différentes.
\end{questions}
\end{corrige}\end{blocLivre}
\begin{blocLivre}{cor-15-4}\begin{corrige}{4}{Lire un motif}
\schema[corrige]{signal-periodique}
\begin{questions}
\item Deux maxima successifs sont séparés de 2,0 ms.
\item $f=1/(2,0\times10^{-3})=\qty{500}{Hz}$.
\end{questions}
\end{corrige}\end{blocLivre}
\begin{blocLivre}{cor-15-5}\begin{corrige}{5}{Mesurer plusieurs périodes}
\begin{questions}
\item $T=12,5/5=\qty{2,50}{ms}$ et $f=\qty{400}{Hz}$.
\item L’incertitude de position des extrémités est répartie sur plusieurs périodes au lieu d’une seule.
\end{questions}
\end{corrige}\end{blocLivre}
\begin{blocLivre}{cor-15-6}\begin{corrige}{6}{Deux instruments}
\schema[corrige]{timbres}
\begin{questions}
\item Même période, donc même fréquence et même hauteur perçue si les deux sont audibles.
\item Leurs formes diffèrent : les timbres peuvent être différents.
\end{questions}
\end{corrige}\end{blocLivre}
\begin{blocLivre}{cor-15-7}\begin{corrige}{7}{Plus aigu}
\begin{questions}
\item Le son de 660 Hz est plus aigu.
\item $T_{220}/T_{660}=660/220=3$. Le son grave possède la période la plus grande.
\end{questions}
\end{corrige}\end{blocLivre}
\begin{blocLivre}{cor-15-8}\begin{corrige}{8}{Une échelle en décibels}
\begin{questions}
\item Elle est multipliée par 100 d’après le tableau.
\item Le niveau est une grandeur sur une échelle non linéaire ; le rapport 70/50 ne donne pas le rapport des intensités.
\end{questions}
\end{corrige}\end{blocLivre}
\clearpage
\begin{blocLivre}{cor-15-9}\begin{corrige}{9}{Une confusion de paramètres}
\begin{questions}
\item Non si la période du motif reste identique.
\item Le son peut être perçu plus fort ; sa hauteur reste liée à la même fréquence.
\end{questions}
\end{corrige}\end{blocLivre}
\begin{blocLivre}{cor-15-10}\begin{corrige}{10}{Un générateur programmé}
\begin{questions}
\item $f=1/0,00125=\qty{800}{Hz}$.
\item Il faut diviser la période par deux : $T=\qty{0,625}{ms}$.
\end{questions}
\end{corrige}\end{blocLivre}
\begin{blocLivre}{cor-15-11}\begin{corrige}{11}{Un son non périodique}
\begin{questions}
\item Non : deux bosses ne définissent pas nécessairement un motif répété.
\item La relation $f=1/T$ s’applique ici à un signal périodique ; sans période identifiable, cette lecture serait arbitraire.
\end{questions}
\end{corrige}\end{blocLivre}
\begin{blocLivre}{cor-15-12}\begin{corrige}{12}{Réduire une exposition}
\demarche{Relier l’exposition au niveau reçu, à la durée et à la distance.}
\begin{questions}
\item S’éloigner, réduire la durée d’exposition et employer une protection adaptée ; diminuer le niveau à la source est également utile.
\item Le niveau reçu et la durée comptent ; le style ne renseigne pas directement sur ces grandeurs.
\end{questions}
\retiens{La hauteur ou le style d’un son ne donne pas directement son niveau d’intensité sonore.}
\end{corrige}\end{blocLivre}

